\chapter{Floating-Point Arithmetic, Conditioning, and Stability}
\label{ch:fp}

Supplementary reading:
\begin{itemize}
  \item \cite{TB} Lectures 12, 13, and 14.
  \item \cite{Demmel} Sections 1.3--1.5.
  \item \cite{GVL} Section 2.4.
\end{itemize}

\section{Why Exact Linear Algebra Is Not Enough}

In a linear algebra course, $\R$ is a field: addition is associative, multiplication distributes over addition, and $x$ and $(x + a) - a$ are the same number. Every derivation in that course relies on these facts.

A computer, however, has no $\R$. Its native objects are bits, $\B \doteq \{0, 1\}$, and any number system it offers must be built from finitely many of them. Approximating an uncountable set by a finite one causes two kinds of trouble:
\begin{enumerate}
  \item \emph{Representation error}: the number $x$ itself cannot be stored, so we store a nearby $\hat{x}$ instead.
  \item \emph{Propagated error}: even if the inputs are exact, the result of every operation has to be pushed back into the finite set.
\end{enumerate}
The second kind has a consequence that is worse than it first appears: \emph{floating-point addition is not associative}. Two algorithms that are mathematically equivalent can therefore return completely different answers on a computer. The question ``is this algorithm correct?'' splits off from the question ``is this algorithm numerically good?'', and this lecture builds the language for the second one. The \emph{condition number} measures how sensitive the problem itself is, and \emph{stability} measures whether an algorithm amplifies errors beyond that.

The topics of the semester are: numerical representation of vectors and matrices; solving linear systems by Gaussian elimination; solving least squares problems by orthogonalization; eigendecompositions; low-rank approximation; iterative methods; and structured matrices.

\section{Vectors and Matrices}

We write $\K$ for either $\R$ or $\C$, so that every result only needs to be stated once. A \emph{vector} $\vec{v} \in \K^m$ is an $m$-tuple of elements of $\K$, representing a position or direction in an $m$-dimensional coordinate system. A \emph{matrix} $A \in \K^{m \times n}$ is an $m \times n$ array of elements of $\K$, which may be thought of as $n$ positions in an $m$-dimensional coordinate system. We identify $\K^{m \times 1} \cong \K^m$, so a column vector is a matrix with one column:
\begin{equation}
  A = \begin{bmatrix} A_{11} & A_{12} \\ A_{21} & A_{22} \\ A_{31} & A_{32} \end{bmatrix} = \begin{bmatrix} \vec{a}_1 & \vec{a}_2 \end{bmatrix}.
\end{equation}
Scalar multiplication and addition act entrywise. For $A \in \K^{m \times n}$, $\vec{v} \in \K^n$, and $B \in \K^{n \times \ell}$,
\begin{equation}
  (A\vec{v})_i = \sum_{k} A_{ik} v_k, \qquad (AB)_{ij} = \sum_k A_{ik} B_{kj},
\end{equation}
so $A$ is a linear map $\K^n \to \K^m$ and $\K^{n \times \ell} \to \K^{m \times \ell}$: it carries a region, together with the vectors $\vec{x}$ in it, to the image vectors $A\vec{x}$.

The entrywise formula is rarely the most useful way to think. Numerical linear algebra relies much more on the \emph{column view}:
\begin{equation}
  A\vec{v} = \sum_{j=1}^{n} v_j \vec{a}_j,
\end{equation}
that is, $A\vec{v}$ is the linear combination of the columns of $A$ with coefficients given by the entries of $\vec{v}$. This view appears throughout the course. The column space $\range(A)$ decides whether $A\vec{x} = \vec{b}$ is solvable, the QR factorization orthogonalizes the columns, and low-rank approximation asks how few directions suffice to approximately span them. Likewise, the $j$-th column of $AB$ is $A$ applied to the $j$-th column of $B$, so a matrix--matrix product is $\ell$ matrix--vector products done at once.

\section{Representing Real Numbers}

We build up to floating-point numbers in four steps, each fixing a shortcoming of the previous one.

\subsection{Unsigned and Signed Integers}

A string $d \in \B^b$ represents the integers $0$ through $2^b - 1$ via
\begin{equation}
  \operatorname{UInt}(d) \doteq \sum_{i=1}^{b} d_i 2^{i-1}.
\end{equation}
To represent negative numbers, we use an offset and give the highest bit a negative weight:
\begin{equation}
  \operatorname{Int}(d) \doteq \sum_{i=1}^{b-1} d_i 2^{i-1} - d_b 2^{b-1},
\end{equation}
which covers $-2^{b-1}$ through $2^{b-1} - 1$. This is \emph{two's complement}. Nonnegative numbers keep their unsigned encoding and the adder circuit is shared with unsigned integers. The price is an asymmetric range: the most negative number, such as \texttt{typemin(Int8)} $= -128$, has no positive counterpart, so $-x$ can overflow back to $x$. (A different kind of offset, the \emph{biased} encoding that stores $x + K$, is used for floating-point exponents; see \Cref{sec:ieee}.)

\subsection{Fixed-Point Numbers}

Fixing the position of the binary point and using $f = \lfloor b/c \rfloor$ of the bits as fractional digits gives
\begin{equation}
  \operatorname{Fixed}(d) \doteq \operatorname{Int}(d) \cdot 2^{-f}.
\end{equation}
Fixed-point numbers have \emph{constant absolute spacing} $2^{-f}$ over the whole range. To represent numbers of size $s$ and $s^{-1}$ at the same time, one needs about $\log_2(s)$ bits, so precision and dynamic range are tied together and the cost grows \emph{linearly} in the dynamic range. Physical simulations routinely span a dozen orders of magnitude, which quickly breaks fixed-point arithmetic.

\subsection{Floating-Point Numbers}

The idea of floating point is to store the order of magnitude separately as an exponent. Given $(p, e) \in \B^{b_p} \times \B^{b_e}$, define
\begin{equation}
  \operatorname{Float}(p, e) \doteq (-1)^{p_1} \left(1 + \sum_{i=2}^{b_p} p_i 2^{-(i-1)}\right) 2^{\operatorname{Int}(e)}.
\end{equation}
The three parts are the sign bit $p_1$, the \emph{significand} (or \emph{mantissa}) in $[1, 2)$, and the power of two. The leading $1$ of the significand is \emph{implicit} and takes no storage. The decimal analogue is $\pm 1.123 \cdot 10^{-4}$.

The key benefit is that covering the range from $s^{-1}$ to $s$ only needs about $\log_2(\log_2(s))$ exponent bits. Fixed point pays a linear price; floating point pays a doubly logarithmic one. This is why floating point dominates scientific computing. A complex number $z \in \C$ is stored as two floating-point numbers. For details, the lecture recommends \cite{Overton}.

\section{The Geometry of Floating-Point Numbers}

Plotting all floating-point numbers on the real line shows that they are \emph{not equally spaced}. Within each \emph{binade} $[2^E, 2^{E+1})$ they are equally spaced, but crossing a power of two doubles the spacing: if the spacing near $1$ is $\eps \doteq 2^{-(b_p - 1)}$, it is $2\eps$ near $2$ and $4\eps$ near $4$.

\begin{center}
\begin{tikzpicture}[x=1.6cm]
  \draw[-{Latex}] (0.8,0) -- (8.6,0);
  \foreach \x in {1,1.25,...,2} \draw (\x,0.12) -- (\x,-0.12);
  \foreach \x in {2,2.5,...,4} \draw (\x,0.12) -- (\x,-0.12);
  \foreach \x in {4,5,...,8} \draw (\x,0.12) -- (\x,-0.12);
  \node[below] at (1,-0.15) {$1$};
  \node[below] at (2,-0.15) {$2$};
  \node[below] at (4,-0.15) {$4$};
  \node[below] at (8,-0.15) {$8$};
  \draw[decorate,decoration={brace,amplitude=4pt}] (1,0.2) -- (1.25,0.2) node[midway,above=4pt] {$\eps$};
  \draw[decorate,decoration={brace,amplitude=4pt}] (2,0.2) -- (2.5,0.2) node[midway,above=4pt] {$2\eps$};
  \draw[decorate,decoration={brace,amplitude=4pt}] (4,0.2) -- (5,0.2) node[midway,above=4pt] {$4\eps$};
\end{tikzpicture}
\end{center}

So although the absolute spacing changes, the \emph{relative spacing is nearly constant}: within one binade it lies between $\eps/2$ and $\eps$. This determines the philosophy behind floating point: it aims for \emph{relative} accuracy, on the premise that every quantity carries an uncertainty proportional to its own size.

\begin{example}[Floating-Point Addition Is Not Associative]
  \label{ex:nonassoc}
  Let $\eps$ be the spacing at $1$ and suppose ties are rounded up. Evaluate $(1 + \eps) + 1 + 2$ in two ways.
  \begin{itemize}
    \item Left to right: $(1 + \eps) + 1 = 2 + \eps$ lies exactly halfway between $2$ and $2 + 2\eps$ and rounds up to $2 + 2\eps$. Adding $2$ gives $4 + 2\eps$, halfway between $4$ and $4 + 4\eps$, which rounds up to $4 + 4\eps$.
    \item Last two terms first: $1 + 2 = 3$ exactly, and $(1 + \eps) + 3 = 4 + \eps$ rounds down to $4$.
  \end{itemize}
  The same expression yields $4$ or $4 + 4\eps$. Multiplication is not associative either, and the distributive law fails as well; commutativity of addition and of multiplication does hold.
\end{example}

The practical consequences are that a parallel reduction depends on the number of threads, that compiler flags such as \texttt{-ffast-math} may reorder sums and change the result, and that bitwise reproducibility across machines takes deliberate effort.

\section{Cancellation}

\begin{example}[Cancellation Exposes Existing Error]
  \label{ex:cancel}
  Suppose $10^6 + 1$ already carries a relative error $\eps$, and we subtract $10^6$:
  \begin{equation}
    (10^6 + 1)(1 + \eps) - 10^6 = 1 + 10^6 \eps.
  \end{equation}
  The exact answer is $1$ and the error is $10^6\eps$. The absolute error is unchanged, but the result is six orders of magnitude smaller than the operands, so the relative error has grown by a factor of $10^6$.
\end{example}

The subtraction itself is innocent: it introduces only $O(\eps)$ new error, and when the operands are close it introduces none at all (see Sterbenz's lemma in \Cref{sec:fp-further}). Cancellation merely \emph{exposes} error that was already present. A more accurate version of the slogan ``avoid cancellation'' is therefore: do not compute a small result as the difference of two large numbers.

\section{Conditioning}
\label{sec:conditioning}

How do errors \emph{propagate}? Consider computing $y = f(x)$ for $f : \R \to \R$. Since the input is rounded before anything else happens, the best we can hope for is $\hat{y} = f(\hat{x})$ with $\abs{\hat{x} - x}/\abs{x} < \eps$. Comparing relative errors,
\begin{equation}
  \frac{\abs{\hat{y} - y}/\abs{y}}{\abs{\hat{x} - x}/\abs{x}} = \frac{\abs{f(\hat{x}) - f(x)}}{\abs{\hat{x} - x}} \cdot \frac{\abs{x}}{\abs{f(x)}} \xrightarrow[\eps \to 0]{} \frac{\abs{x}\abs{f'(x)}}{\abs{f(x)}}.
\end{equation}

\begin{definition}[Condition Number]
  \label{def:cond}
  The (relative) \emph{condition number} of $f$ at $x$ is
  \begin{equation}
    \kappa_f(x) \doteq \frac{\abs{x}\abs{f'(x)}}{\abs{f(x)}}.
  \end{equation}
  The problem is \emph{well-conditioned} at $x$ if $\kappa_f(x) \approx 1$ and \emph{ill-conditioned} if $\kappa_f(x) \gg 1$.
\end{definition}

The condition number depends only on the problem and the input, not on the algorithm. This is the most important point of the lecture: no algorithm can rescue an ill-conditioned problem, because the difficulty is intrinsic.

\begin{example}[Conditioning of Subtraction]
  Let $f(x) = x - 10^6$ at $x = 10^6 + 1$. Then
  \begin{equation}
    \kappa_f(10^6 + 1) = \frac{\abs{10^6 + 1} \cdot 1}{1} \approx 10^6 \gg 1.
  \end{equation}
  Cancellation is thus a conditioning phenomenon, not a defect of floating point. Even an algorithm with infinite precision turns a relative uncertainty of $10^{-16}$ in $x$ into a relative uncertainty of $10^{-10}$ in the output.
\end{example}

\section{Stability}
\label{sec:stability}

Conditioning is about the problem; stability is about the algorithm. Let $\hat{f}$ be the floating-point algorithm we actually run: it acts on $\hat{x}$ and rounds at every internal step.

\begin{definition}[Stability]
  An algorithm $\hat{f}$ for $f$ is \emph{stable} if
  \begin{equation}
    \frac{\abs{\hat{f}(\hat{x}) - y}/\abs{y}}{\abs{\hat{x} - x}/\abs{x}} \lesssim \kappa_f(x),
  \end{equation}
  that is, if it does not amplify input perturbations more than the problem itself does. Otherwise $\hat{f}$ is \emph{unstable}.
\end{definition}

\begin{example}[An Unstable Algorithm for the Identity]
  \label{ex:unstable-identity}
  Let $f(x) = x$ and $\hat{f}(x) = (x + 10^6) - 10^6$. Mathematically $\hat{f} = f$, and the identity has $\kappa_f \equiv 1$, the best possible conditioning. But $\hat{f}$ first adds $x$ to $10^6$, which pushes the low-order bits of $x$ out of the significand for good, and then subtracts $10^6$ again. For $x = 1$ the relative error is $O(10^6 \eps)$ instead of $O(\eps)$.

  The diagnosis: \emph{the algorithm splits a well-conditioned problem into ill-conditioned intermediate steps}. Writing $\hat{f} = g_2 \circ g_1$ with $g_1(x) = x + 10^6$ and $g_2(z) = z - 10^6$, the step $g_1$ is well-conditioned but $g_2$ has condition number about $10^6$ at $z \approx 10^6$. A small overall condition number does not mean every step is well-conditioned, and the algorithm's error accumulates step by step. This view in terms of intermediate condition numbers will return when we compare Gram--Schmidt with the normal equations.
\end{example}

``Unstable'' and ``ill-conditioned'' are different diagnoses with different remedies: the former calls for a different algorithm, the latter for a different formulation of the problem (reparametrization, regularization, or higher precision).

\section{(SUPPLEMENT) IEEE 754 Formats}
\label{sec:ieee}

The formats in common use are listed below; $b_p$ counts the implicit bit, so a significand field with $b_p - 1$ stored bits provides $b_p$ bits of precision.
\begin{center}
\begin{tabular}{lcccccc}
  \toprule
  Format & Total bits & Exponent bits & Stored significand bits & $b_p$ & Bias & Decimal digits \\
  \midrule
  Float64 & 64 & 11 & 52 & 53 & 1023 & $\approx 15.9$ \\
  Float32 & 32 & 8 & 23 & 24 & 127 & $\approx 7.2$ \\
  Float16 & 16 & 5 & 10 & 11 & 15 & $\approx 3.3$ \\
  BFloat16 & 16 & 8 & 7 & 8 & 127 & $\approx 2.4$ \\
  \bottomrule
\end{tabular}
\end{center}
The exponent field uses a biased encoding, and its two extreme values are reserved:
\begin{itemize}
  \item An all-zeros exponent encodes $\pm 0$ (if the significand is also zero) or a \emph{subnormal} number, which drops the implicit leading $1$ and uses $0.f \times 2^{E_{\min}}$. Subnormals provide \emph{gradual underflow}, which guarantees $x \ne y \implies x - y \ne 0$.
  \item An all-ones exponent encodes $\pm\infty$ (zero significand) or \texttt{NaN} (nonzero significand).
\end{itemize}
The normal range of Float64 is therefore roughly $2.2 \times 10^{-308}$ to $1.8 \times 10^{308}$. IEEE defines five exceptions (invalid, division by zero, overflow, underflow, inexact), and specifies that $\infty$ and \texttt{NaN} propagate through arithmetic instead of halting it. Demmel~\cite[\S1.5]{Demmel} stresses the value of this: an algorithm can run to completion and check its result at the end, instead of branching at every step. Note that \texttt{NaN != NaN} is always true, so one must test with \texttt{isnan}.

\section{(SUPPLEMENT) Machine Epsilon and the Model of Arithmetic}

\subsection{Two Conventions for Machine Epsilon}

Two different quantities are both called ``machine epsilon'' in the literature, and they differ by a factor of $2$.

\begin{definition}[Machine Epsilon and Unit Roundoff]
  \label{def:eps}
  Let $b_p$ be the precision of a floating-point format.
  \begin{itemize}
    \item The \emph{machine epsilon} $\eps = \epsm \doteq 2^{-(b_p - 1)}$ is the distance from $1$ to the next larger floating-point number. For Float64, $\eps = 2^{-52} \approx 2.22 \times 10^{-16}$. This is the convention of the lecture, of Julia's \texttt{eps(Float64)}, and of C's \texttt{DBL\_EPSILON}.
    \item The \emph{unit roundoff} $u \doteq 2^{-b_p} = \eps/2$ is the largest relative rounding error, i.e., half a spacing. For Float64, $u = 2^{-53} \approx 1.11 \times 10^{-16}$. Trefethen and Bau~\cite{TB} call this quantity $\epsilon_{\text{machine}}$, and Golub and Van Loan~\cite{GVL} call it the unit roundoff.
  \end{itemize}
\end{definition}

Any statement of the form $O(\eps)$ is unaffected by the choice, but constants copied from different textbooks must be read with care.

The IEEE default rounding mode is \emph{round to nearest, ties to even}, under which both orders in \Cref{ex:nonassoc} give $4$. Non-associativity does not depend on the rounding mode, however. With ties to even,
\begin{equation}
  (1 + u) + u = 1 \qquad \text{but} \qquad 1 + (u + u) = 1 + \eps,
\end{equation}
since on the left each intermediate result is a tie that rounds back to $1$, while on the right $u + u = \eps$ and $1 + \eps$ are both exact.

\subsection{The Fundamental Axiom}

Error analysis needs rounding to be captured by a rule we can reason with. Let $\fl(x)$ denote the floating-point number nearest to $x$. Because the relative spacing is bounded, every $x$ in the normal range satisfies
\begin{equation}
  \fl(x) = x(1 + \delta), \qquad \abs{\delta} \le u. \label{eq:fl}
\end{equation}
IEEE requires $+$, $-$, $\times$, $\div$, and $\sqrt{\cdot}$ to be \emph{correctly rounded}: the exact result is computed in infinite precision and then rounded once. This gives the following model, stated in the same form in \cite[Lecture 13]{TB}, \cite[\S2.4]{GVL}, and \cite[\S1.5]{Demmel}.

\begin{theorem}[Fundamental Axiom of Floating-Point Arithmetic]
  \label{thm:axiom}
  For all floating-point numbers $x, y$ and $\circledast \in \{+, -, \times, \div\}$,
  \begin{equation}
    \fl(x \circledast y) = (x \circledast y)(1 + \delta), \qquad \abs{\delta} \le u.
  \end{equation}
\end{theorem}

The axiom has an immediate \emph{backward} reading. For addition,
\begin{equation}
  \fl(x + y) = (x + y)(1 + \delta) = x(1 + \delta) + y(1 + \delta),
\end{equation}
so the computed sum is \emph{exactly} the sum of two slightly perturbed inputs. Pushing errors back onto the inputs in this way is the main thread of error analysis in this course; see \Cref{sec:backward}.

\subsection{Accumulation of Errors}

A single operation incurs an error of size $O(u)$. For $n$ operations, Golub and Van Loan write
\begin{equation}
  \gamma_n \doteq \frac{nu}{1 - nu} \approx nu \qquad (nu \ll 1).
\end{equation}

\begin{proposition}[Error in Inner Products]
  \label{prop:dot}
  For $\vec{x}, \vec{y} \in \R^n$,
  \begin{equation}
    \abs{\fl(\vec{x}^\T \vec{y}) - \vec{x}^\T \vec{y}} \le \gamma_n \sum_{i=1}^{n} \abs{x_i}\abs{y_i} = \gamma_n \abs{\vec{x}}^\T \abs{\vec{y}}.
  \end{equation}
  Similarly, for matrices, $\abs{\fl(AB) - AB} \le \gamma_n \abs{A}\abs{B}$ entrywise.
\end{proposition}

The \emph{shape} of this bound matters much more than its constant. The right-hand side involves $\abs{\vec{x}}^\T\abs{\vec{y}}$, not $\abs{\vec{x}^\T\vec{y}}$. If all terms have the same sign, the two agree and the relative error is $O(nu)$, which is fine. If instead $\vec{x}^\T\vec{y}$ is a small number obtained by cancelling many terms of both signs, then $\abs{\vec{x}}^\T\abs{\vec{y}} \gg \abs{\vec{x}^\T\vec{y}}$ and the relative error can be arbitrarily bad. The linear factor $n$ is almost never the real issue (it is a very pessimistic worst case, and random errors grow more like $\sqrt{n}$). The real issue is always \emph{cancellation}.

\section{(SUPPLEMENT) Backward Stability}
\label{sec:backward}

The textbooks~\cite[Lecture 14]{TB} mostly use a stronger notion than the stability of \Cref{sec:stability}.

\begin{definition}[Backward Stability]
  \label{def:backward-stable}
  An algorithm $\hat{f}$ for $f$ is \emph{backward stable} if for every $x$ there exists $\tilde{x}$ with
  \begin{equation}
    \hat{f}(x) = f(\tilde{x}) \qquad \text{and} \qquad \frac{\abs{\tilde{x} - x}}{\abs{x}} = O(u).
  \end{equation}
\end{definition}

In words: \emph{the computed answer is the exact answer to a nearby problem}. The backward reading of \Cref{thm:axiom} is the simplest example. Backward stability is the gold standard of numerical linear algebra because it cleanly separates the problem from the algorithm, and it yields the classical rule of thumb
\begin{equation}
  \text{forward relative error} \lesssim \kappa_f(x) \cdot u,
\end{equation}
or, in a form that is easier to remember,
\begin{equation}
  \text{accuracy} \approx \text{condition number (problem)} \times \text{backward error (algorithm)}.
\end{equation}
A backward stable algorithm still returns inaccurate answers for ill-conditioned problems, and this is \emph{not the algorithm's fault}; it cannot be improved upon. Conversely, when a result is inaccurate, the first step should be to estimate the condition number, not to suspect the code.

\begin{center}
\begin{tabular}{l p{0.2\textwidth} p{0.34\textwidth} l}
  \toprule
  Concept & Belongs to & Question it answers & Measure \\
  \midrule
  Conditioning & problem $f$ and input $x$ & how much input perturbations are intrinsically amplified & $\kappa_f(x)$ \\
  Stability & algorithm $\hat{f}$ & whether the implementation amplifies errors further & backward error $O(u)$ \\
  Accuracy & final result & whether the computed number is right & $\lesssim \kappa \cdot u$ \\
  \bottomrule
\end{tabular}
\end{center}

Some common condition numbers:
\begin{center}
\begin{tabular}{lll}
  \toprule
  Problem & $\kappa$ & Remark \\
  \midrule
  $f(x) = \sqrt{x}$ & $1/2$ & always well-conditioned; even shrinks errors \\
  $f(x) = x^n$ & $n$ & grows linearly with the power \\
  $f(x) = e^x$ & $\abs{x}$ & ill-conditioned for large $\abs{x}$ \\
  $f(x_1, x_2) = x_1 - x_2$ & $\dfrac{\abs{x_1} + \abs{x_2}}{\abs{x_1 - x_2}}$ & cancellation (componentwise perturbations) \\
  $\vec{x} \mapsto A\vec{x}$ & $\dfrac{\norm{A}\norm{\vec{x}}}{\norm{A\vec{x}}} \le \kappa(A)$ & depends on the direction of $\vec{x}$ \\
  solve $A\vec{x} = \vec{b}$ & $\kappa(A) = \norm{A}\norm{A^{-1}}$ & \\
  polynomial roots & can be enormous & Wilkinson's polynomial \\
  \bottomrule
\end{tabular}
\end{center}
For $f : \R^n \to \R^m$ one replaces $f'$ by the Jacobian $J(\vec{x})$: the absolute condition number is $\hat{\kappa} = \norm{J(\vec{x})}$ and the relative condition number is $\kappa = \norm{J(\vec{x})}\norm{\vec{x}}/\norm{f(\vec{x})}$~\cite[Lecture 12]{TB}.

\section{(SUPPLEMENT) Further Topics}
\label{sec:fp-further}

\begin{strategy}[Rewriting to Avoid Cancellation]
  Common traps and their fixes:
  \begin{itemize}
    \item The root $\frac{-b + \sqrt{b^2 - 4ac}}{2a}$ of a quadratic cancels when $b^2 \gg 4ac$; compute the other root first and recover this one from Vieta's relation $x_1 x_2 = c/a$.
    \item The variance formula ``mean of squares minus square of mean'' cancels; use a two-pass algorithm or Welford's recurrence.
    \item $1 - \cos x$, $\sqrt{x + 1} - \sqrt{x}$, and $\log(1 + x)$ cancel as $x \to 0$; use a half-angle identity, rationalize, or call \texttt{log1p} and \texttt{expm1}.
  \end{itemize}
\end{strategy}

\begin{lemma}[Sterbenz]
  If $x, y$ are floating-point numbers with $y/2 \le x \le 2y$, then $x - y$ is a floating-point number, so $\fl(x - y) = x - y$ exactly.
\end{lemma}
This confirms again that cancellation does not lose new precision; it only exposes old errors.

\paragraph{Useful techniques.}
\begin{itemize}
  \item \emph{Kahan compensated summation} carries an extra variable that tracks the low-order bits lost in each rounding, reducing the error bound for a sum from $O(nu)$ to $O(u)$, independent of $n$. It costs four times as many operations, and \texttt{-ffast-math} may optimize it away.
  \item \emph{Fused multiply-add} (FMA) computes $a \cdot b + c$ with one rounding instead of two. It is a basic instruction on modern CPUs and GPUs, and it is why two mathematically equivalent expressions, or even \texttt{a*b - a*b}, may give different results.
  \item \texttt{-ffast-math} (or Julia's \texttt{@fastmath}) lets the compiler pretend that floating point is associative and distributive, trading correctness for speed. It should usually stay off in numerical code.
\end{itemize}

\paragraph{Low-precision formats.} Float16 and BFloat16 both use $16$ bits, but BFloat16 spends them on the exponent: it has $8$ exponent bits, the same range as Float32, and only $8$ bits of precision. This is the right trade-off for deep learning, where the first failure in training is usually overflow or underflow of gradients (a range problem) rather than lack of precision. Float16 has only $5$ exponent bits and a largest value of about $65504$, so training in Float16 requires loss scaling while BFloat16 does not. The accumulation bound $\gamma_n\abs{A}\abs{B}$ of \Cref{prop:dot} explains another convention, \emph{multiply in low precision, accumulate in high precision}: the factor $n$ comes from the length of the summation chain, and reduction dimensions of thousands are common, so tensor cores take BFloat16 or FP8 inputs but accumulate in Float32.

\paragraph{Experiments in Julia.}
\begin{lstlisting}
bitstring(1.0)                    # sign, exponent, and significand fields
eps(Float64), eps(1.0), eps(2.0)  # last two differ by 2: spacing doubles
nextfloat(1.0) - 1.0              # == eps(Float64) == 2^-52
0.1 + 0.2 == 0.3                  # false
(1.0 + eps()/2) + eps()/2 == 1.0 + (eps()/2 + eps()/2)   # false: not associative
1e16 + 1.0 - 1e16                 # compare with 1e6 + 1.0 - 1e6
Float16(65504), Float16(65520)    # the second overflows to Inf
nextfloat(Float64(0))             # smallest subnormal, about 5e-324
\end{lstlisting}
